% \iffalse meta-comment
%
% 四张轴承图的 TikZ 源码。坐标以原 PNG 左上角为原点，单位对应原图像素；
% 原图尺寸分别为 600×274、576×276、582×311、621×319。
% 轮廓、标注中心与字号按原图校准，输出以 1 bp 对应 1 px 保留纵横比。
% make figures 编译为 assets/ 下的 PDF；字体由 ctex 自动选择，不固定 fontset。
% 涡室截面只提取 tikzpicture，分发到示例 figures/volute-section.tex 供正文 input。
%
% \fi
% \changes{v4.0a}{2026/10/07}{四张轴承示例图改用 TikZ 源码生成矢量 PDF，并示范正文内联绘图；涡室截面以独立 TikZ 文件供正文引用；移除高尔夫插图、支付宝二维码及其 EPS 资源}
%    \begin{macrocode}
%<*figure-common>
\documentclass[tikz,border=0pt]{standalone}
\usepackage{ctex}
\usetikzlibrary{arrows.meta,patterns.meta}
\tikzset{
  original/.style={x=1bp,y=-1bp,line width=1.5bp,line cap=butt,line join=miter},
  wall/.style={draw,pattern={Lines[angle=45,distance=20bp,line width=1bp]}},
  lower wall/.style={wall,pattern={Lines[angle=-45,distance=20bp,line width=1bp]}},
  restrictor/.style={draw,pattern={Lines[angle=-65,distance=8bp,line width=1bp]}},
  flow/.style={-{Triangle[length=11bp,width=7bp]},line width=1.6bp},
  leader/.style={-{Triangle[length=12bp,width=8bp]},dash pattern=on 3bp off 3bp,line width=1bp},
  every node/.style={font=\songti\fontsize{28bp}{28bp}\selectfont,inner sep=0pt,outer sep=0pt},
}
\begin{document}
%</figure-common>
%<*aerostatic>
\begin{tikzpicture}[original]
  \path[use as bounding box] (0,0) rectangle (600,274);
  \begin{pgfinterruptboundingbox}
    \draw[wall] (188,110) -- (370,110) -- (370,194) -- (188,194) -- cycle;
    \draw[wall] (412,110) -- (590,110) -- (590,194) -- (412,194) -- cycle;
    \draw (370,110) -- (412,110);
    \draw[restrictor] (370,166) rectangle (412,189);
    \draw (370,194) -- (412,194);
    \draw[lower wall] (188,218) rectangle (590,260);
    \draw[flow] (379,0) -- (379,82);
    \draw[flow] (400,0) -- (400,82);
    \node at (494,23) {加压气体};
    \node at (234,23) {节流器};
    \draw (300,40) -- (328,94);
    \draw[leader] (328,94) -- (380,179);
    \draw[flow] (332,205) -- (231,205);
    \draw[flow] (525,205) -- (542,205);
    \node[anchor=west] at (5,208) {气体流动方向};
  \end{pgfinterruptboundingbox}
\end{tikzpicture}
%</aerostatic>
%<*aerodynamic>
\begin{tikzpicture}[original]
  \path[use as bounding box] (0,0) rectangle (576,276);
  \begin{pgfinterruptboundingbox}
    \draw[wall] (154,53) -- (328,53) -- (328,109) -- (154,95) -- cycle;
    \draw[wall] (372,53) -- (564,53) -- (564,125) -- (372,113) -- cycle;
    \draw (328,53) -- (372,53);
    \draw[restrictor] (328,80) rectangle (372,101);
    \draw (328,109) -- (372,113);
    \fill[pattern={Lines[angle=-45,distance=20bp,line width=1bp]}] (154,165) rectangle (564,206);
    \draw (154,206) -- (154,165) -- (564,165) -- (564,206);
    \node at (281,23) {节流器};
    \draw[leader] (314,45) -- (339,97);
    \node at (510,23) {楔状间隙};
    \draw[leader] (499,39) -- (417,135);
    \node[anchor=west] at (10,139) {气体流动方向};
    % 原位图的细线有采样断裂；流向统一用完整实线箭头表示。
    \draw[flow] (192,113) -- (357,123);
    \draw[flow] (192,132) -- (357,140);
    \draw[flow] (192,155) -- (358,155);
    \draw[flow] (422,221) -- (441,221);
    \node at (356,260) {移动方向};
  \end{pgfinterruptboundingbox}
\end{tikzpicture}
%</aerodynamic>
%<*hybrid>
\begin{tikzpicture}[original]
  \path[use as bounding box] (0,0) rectangle (582,311);
  \begin{pgfinterruptboundingbox}
    \draw[wall] (154,56) -- (329,56) -- (329,115) -- (154,101) -- cycle;
    \draw[wall] (373,56) -- (566,56) -- (566,137) -- (373,120) -- cycle;
    \draw (329,56) -- (373,56);
    \draw[restrictor] (329,86) rectangle (373,108);
    \draw (329,115) -- (373,120);
    \draw[lower wall] (154,184) rectangle (566,230);
    \node at (281,21) {节流器};
    \draw[leader] (313,48) -- (338,104);
    \node at (510,21) {楔状间隙};
    \draw[leader] (500,40) -- (418,148);
    \node[anchor=west] at (10,151) {气体流动方向};
    \draw[flow] (192,123) -- (358,135);
    \draw[flow] (192,145) -- (358,152);
    \draw[flow] (192,170) -- (358,170);
    \draw[flow] (423,245) -- (442,245);
    \node at (356,291) {移动方向};
  \end{pgfinterruptboundingbox}
\end{tikzpicture}
%</hybrid>
%<*squeeze-film>
\begin{tikzpicture}[original]
  \path[use as bounding box] (0,0) rectangle (621,319);
  \begin{pgfinterruptboundingbox}
    \draw[wall] (192,128) rectangle (605,223);
    \draw (192,111) rectangle (605,128);
    \draw (192,223) rectangle (605,242);
    \draw[{Triangle[length=13bp,width=9bp]}-{Triangle[length=13bp,width=9bp]},line width=1.8bp]
      (409,14) -- (409,83);
    \node at (506,32) {振动方向};
    \draw[flow] (338,255) -- (234,255);
    \draw[flow] (453,255) -- (557,255);
    \node[anchor=west] at (13,264) {气体流动方向};
    \begin{scope}
      \clip (191,271) rectangle (603,291);
      \foreach \x in {160,190,...,610}
        \draw[line width=1bp] (\x,271) -- ++(40,20);
    \end{scope}
    \draw (191,271) -- (603,271);
  \end{pgfinterruptboundingbox}
\end{tikzpicture}
%</squeeze-film>
% 此段没有文档类和 document 环境，也不依赖前面的绘图样式。
% 调用方在导言区加载 TikZ 的 arrows.meta 库。
%<*volute-section>
\begin{tikzpicture}[
  x=1bp,y=-1bp,line width=1.2bp,
  dimension/.style={{Triangle[length=5bp,width=3bp]}-{Triangle[length=5bp,width=3bp]},line width=.65bp},
  every node/.style={font=\songti\fontsize{16bp}{18bp}\selectfont,inner sep=1bp},
]
  \path[use as bounding box] (0,0) rectangle (753,269);
  \begin{pgfinterruptboundingbox}
    \foreach \i/\cx/\title in {0/112/{(a) 任意形状},1/293/{(b) 矩形},2/473/{(c) 梯形},3/643/{(d) 圆形}} {
      \begin{scope}[shift={(\cx,0)}]
        \draw[line width=.6bp,dash pattern=on 14bp off 5bp on 2bp off 5bp] (0,10) -- (0,224);
        \ifcase\i
          \def\bottom{106}\def\layer{79}\def\halfwidth{25}
          \draw (-32,106) .. controls (-20,80) and (-36,68) .. (-38,48)
            .. controls (-40,9) and (40,7) .. (38,48)
            .. controls (36,68) and (20,81) .. (32,106) -- cycle;
        \or
          \def\bottom{98}\def\layer{79}\def\halfwidth{28}
          \draw (-28,98) -- (-28,47)
            .. controls (-28,10) and (28,10) .. (28,47)
            -- (28,98) -- cycle;
        \or
          \def\bottom{98}\def\layer{79}\def\halfwidth{28}
          \draw (-23,98) -- (-36,57)
            .. controls (-47,8) and (47,8) .. (36,57)
            -- (23,98) -- cycle;
        \or
          \def\bottom{98}\def\layer{74}\def\halfwidth{30}
          \draw (-23,98) -- (-30.31,77.22)
            .. controls (-34.65,71.01) and (-37,63.55) .. (-37,56)
            .. controls (-37,35.57) and (-20.43,19) .. (0,19)
            .. controls (20.43,19) and (37,35.57) .. (37,56)
            .. controls (37,63.55) and (34.65,71.01) .. (30.31,77.22)
            -- (23,98) -- cycle;
        \fi
        \begin{scope}
          \clip (-\halfwidth,\layer) rectangle (\halfwidth,{\layer+6});
          \foreach \x in {-36,-33,...,36}
            \draw[line width=.5bp] (\x,\layer) -- ++(-3.5,6);
        \end{scope}
        \draw[line width=.6bp] (-\halfwidth,\layer) rectangle (\halfwidth,{\layer+6});
        \draw[line width=.6bp] (-78,19) -- (0,19);
        \draw[line width=.6bp] (-48,\layer) -- (-\halfwidth,\layer);
        \draw[line width=.6bp] (-48,{\layer+6}) -- (-\halfwidth,{\layer+6});
        \draw[line width=.6bp] (23,\bottom) -- (48,\bottom);
        \draw[line width=.6bp] (-78,224) -- (49,224);
        \draw[dimension] (-74,19) -- node[rotate=90,above] {$R_c$} (-74,224);
        \draw[dimension] (-46,19) -- node[rotate=90,above] {$\mathrm{d}R$} (-46,\layer);
        \draw[dimension] (-46,{\layer+6}) -- node[rotate=90,above] {$R$} (-46,224);
        \draw[dimension] (44,\bottom) -- node[rotate=90,above] {$R_3$} (44,224);
        \draw[dimension] (-\halfwidth,61) -- node[above] {$b$} (\halfwidth,61);
        \node at (0,250) {\title};
      \end{scope}
    }
  \end{pgfinterruptboundingbox}
\end{tikzpicture}
%</volute-section>
%<*figure-end>
\end{document}
%</figure-end>
%    \end{macrocode}
